% Source: https://artofproblemsolving.com/community/c3h1641956_an_interesting_generalization
% Stable identity: aops-topic-1641956-q01
Let $P=abc\,(a^3-b^3)(b^3-c^3)(c^3-a^3)$. We show that $84=2^2\cdot3\cdot7$ always divides $P$ and that nothing larger does.

\textbf{Divisible by $7$.} If $7$ divides one of $a,b,c$, then $7\mid abc$. Otherwise each cube is $\equiv\pm1\pmod 7$ (the cubes of $1,2,\dots,6$ are $\equiv1,1,6,1,6,6$). Among the three cubes two are congruent, and their difference is divisible by $7$.

\textbf{Divisible by $3$.} If $3$ divides one of $a,b,c$ we are done. Otherwise $a,b,c\in\{1,2\}\pmod 3$, so two of them, say $a,b$, are congruent, and $3\mid a-b\mid a^3-b^3$.

\textbf{Divisible by $4$.} If at least two of $a,b,c$ are even, then $4\mid abc$. If exactly one is even, then $2\mid abc$, and the other two are odd, so the difference of their cubes is even; this gives $4$. If all three are odd, all three differences of cubes are even, giving $8$.

Since $3$, $4$, $7$ are pairwise coprime, $84\mid P$ always.

\textbf{Nothing larger.} For $(a,b,c)=(1,2,3)$, $|P|=6\cdot7\cdot19\cdot26=20748=2^2\cdot3\cdot7\cdot13\cdot19$. For $(a,b,c)=(1,2,4)$, $|P|=8\cdot7\cdot56\cdot63=2^6\cdot3^2\cdot7^3$. The greatest common divisor of these two values is $2^2\cdot3\cdot7$, so the answer is $\boxed{84}$.
